\section{Settlement and Verification}
\label{sec:settlement}

\subsection{Epoch Mechanics}

Time advances in \emph{epochs}: bounded settlement windows in which the off-chain engine matches intents, runs state-transition programs, and commits a snapshot on-chain. Epoch duration is configurable per market, trading latency against proof cost.

Each snapshot carries five Merkle roots: executable transfers, internal market state, debt records, liquidation events, and CDP deployments.\footnote{Internal state uses ZK-friendly Merkle structures; on-chain roots use keccak-hashed binary trees.} Posting a new epoch requires a valid \emph{fact} in the facts registry, attesting correct state transition via zero-knowledge proof.

Execution is inclusion-proof driven: anyone can execute a committed leaf by supplying a Merkle proof against the epoch root stored on-chain.

Each epoch advance processes three concurrent generations: finalizing loans from two epochs prior, settling principal from one epoch prior, and matching new intents. When the solver is the vault operator, all three compress into a single epoch via an atomic multicall.

\subsection{Settlement Programs}

On-chain contracts are settlement executors: they verify Merkle proofs and execute committed transfers and hook chains but contain no market-specific logic. All behavior---interest computation, liquidation detection, oracle selection, collateral policy---lives in zero-knowledge programs whose outputs the contracts verify.

Programs split by economic cadence. \textbf{Matching} is one-time per loan---transfers, debt issuance, CDP authorization, and criteria-compatibility evaluation (Section~\ref{sec:architecture}). \textbf{Interest accrual} is periodic per debt per time window. \textbf{Liquidation} is event-triggered when health breaches produce liquidation leaves. \textbf{Close} is terminal after repayments and accrual catch up. \textbf{Compliance} is periodic per vault epoch, attesting operator behavior against the mandate (Section~\ref{sec:lp}).

All programs submit results through a single \emph{facts registry} that accepts a fact hash and returns whether it is valid. Two fact types carry all market state. A \textbf{state-transition fact} authorizes an epoch advance, binding the settlement program, chain, issuer, and the previous and next epoch roots (each encoding the full quintuple of Merkle roots). An \textbf{interest fact} authorizes accrual on one debt, binding the rate program, chain, debt identifier, accrual window, and computed interest amount. The on-chain contract reconstructs the expected fact from its fixed parameters, checks the registry, and applies values only if valid.

\subsection{Market and Loan Contracts}

Each market has two contracts. The \textbf{collateral manager issuer} (collateral chain) advances epoch state, executes Merkle-authorized transfers, and deploys per-loan collateral managers. The \textbf{debt issuer} (principal chain) advances epoch state and deploys per-loan issued debt tokens. Neither contains market-specific logic.

Each loan also has two contracts. The \textbf{collateral manager} (CDP) holds pledged assets and executes hook-composed pipelines; permitted operations are defined by the market's ZK programs. The \textbf{issued debt} is a per-loan ERC-20 tracking accrual via verified facts and partitioning repayment into principal reserve, interest reserve, and protocol fees (Section~\ref{sec:tokens}).

\subsection{Oracle Freedom and Permissionless Markets}

Oracle logic is embedded in the market's zero-knowledge programs. Programs read price or rate data from on-chain contracts via storage proofs, enforcing freshness inside the proof rather than trusting off-chain feeds. Markets can be created permissionlessly: the protocol verifies program outputs but does not audit program logic. Assessing a market's risk properties is the participant's responsibility.

\subsection{System Invariants}

Three invariants hold across all states:
\begin{enumerate}
  \item \textbf{Token conservation.} For every asset, total balances across all market contracts equal deposited collateral plus supplied principal minus withdrawals. No operation creates or destroys tokens.
  \item \textbf{Epoch monotonicity.} Each market's epoch advances strictly by one, only when the facts registry contains a valid state-transition fact linking the previous and next roots. No epoch can be skipped, replayed, or forked.
  \item \textbf{Position isolation.} Each loan's contracts are independently deployed, independently liquidatable, and share no mutable state with any other loan. One position's insolvency cannot affect another.
\end{enumerate}
